%-----------------------------------------------------------------------------------------------------------------------------------------------%
%	The MIT License (MIT)
%
%	Copyright (c) 2021 Jitin Nair
%
%	Permission is hereby granted, free of charge, to any person obtaining a copy
%	of this software and associated documentation files (the "Software"), to deal
%	in the Software without restriction, including without limitation the rights
%	to use, copy, modify, merge, publish, distribute, sublicense, and/or sell
%	copies of the Software, and to permit persons to whom the Software is
%	furnished to do so, subject to the following conditions:
%	
%	THE SOFTWARE IS PROVIDED "AS IS", WITHOUT WARRANTY OF ANY KIND, EXPRESS OR
%	IMPLIED, INCLUDING BUT NOT LIMITED TO THE WARRANTIES OF MERCHANTABILITY,
%	FITNESS FOR A PARTICULAR PURPOSE AND NONINFRINGEMENT. IN NO EVENT SHALL THE
%	AUTHORS OR COPYRIGHT HOLDERS BE LIABLE FOR ANY CLAIM, DAMAGES OR OTHER
%	LIABILITY, WHETHER IN AN ACTION OF CONTRACT, TORT OR OTHERWISE, ARISING FROM,
%	OUT OF OR IN CONNECTION WITH THE SOFTWARE OR THE USE OR OTHER DEALINGS IN
%	THE SOFTWARE.
%	
%
%-----------------------------------------------------------------------------------------------------------------------------------------------%

%----------------------------------------------------------------------------------------
%	DOCUMENT DEFINITION
%----------------------------------------------------------------------------------------

% article class because we want to fully customize the page and not use a cv template
\documentclass[a4paper,12pt]{article}

%----------------------------------------------------------------------------------------
%	FONT
%----------------------------------------------------------------------------------------

% % fontspec allows you to use TTF/OTF fonts directly
% \usepackage{fontspec}
% \defaultfontfeatures{Ligatures=TeX}

% % modified for ShareLaTeX use
% \setmainfont[
% SmallCapsFont = Fontin-SmallCaps.otf,
% BoldFont = Fontin-Bold.otf,
% ItalicFont = Fontin-Italic.otf
% ]
% {Fontin.otf}

%----------------------------------------------------------------------------------------
%	PACKAGES
%----------------------------------------------------------------------------------------
\usepackage{url}
\usepackage{parskip} 	

%other packages for formatting
\RequirePackage{color}
\RequirePackage{graphicx}
\usepackage[dvipsnames]{xcolor}
\usepackage[scale=0.9]{geometry}

%tabularx environment
\usepackage{tabularx}

%for lists within experience section
\usepackage{enumitem}

% centered version of 'X' col. type
\newcolumntype{C}{>{\centering\arraybackslash}X} 

%to prevent spillover of tabular into next pages
\usepackage{supertabular}
\usepackage{tabularx}
\newlength{\fullcollw}
\setlength{\fullcollw}{0.47\textwidth}

%custom \section
\usepackage{titlesec}				
\usepackage{multicol}
\usepackage{multirow}

%CV Sections inspired by: 
%http://stefano.italians.nl/archives/26
\titleformat{\section}{\Large\scshape\raggedright}{}{0em}{}[\titlerule]
\titlespacing{\section}{0pt}{10pt}{10pt}

%for publications
\usepackage[style=authoryear,sorting=ynt, maxbibnames=2]{biblatex}

%Setup hyperref package, and colours for links
\usepackage[unicode, draft=false]{hyperref}
\definecolor{linkcolour}{rgb}{0,0.2,0.6}
\hypersetup{colorlinks,breaklinks,urlcolor=linkcolour,linkcolor=linkcolour}
\addbibresource{citations.bib}
\setlength\bibitemsep{1em}

%for social icons
\usepackage{fontawesome5}

%debug page outer frames
%\usepackage{showframe}


% job listing environments
\newenvironment{jobshort}[2]
    {
    \begin{tabularx}{\linewidth}{@{}l X r@{}}
    \textbf{#1} & \hfill &  #2 \\[3.75pt]
    \end{tabularx}
    }
    {
    }

\newenvironment{joblong}[2]
    {
    \begin{tabularx}{\linewidth}{@{}l X r@{}}
    \textbf{#1} & \hfill &  #2 \\[3.75pt]
    \end{tabularx}
    \begin{minipage}[t]{\linewidth}
    \begin{itemize}[nosep,after=\strut, leftmargin=1em, itemsep=3pt,label=--]
    }
    {
    \end{itemize}
    \end{minipage}    
    }



%----------------------------------------------------------------------------------------
%	BEGIN DOCUMENT
%----------------------------------------------------------------------------------------
\begin{document}

% non-numbered pages
\pagestyle{empty} 

%----------------------------------------------------------------------------------------
%	TITLE
%----------------------------------------------------------------------------------------

% \begin{tabularx}{\linewidth}{ @{}X X@{} }
% \huge{Your Name}\vspace{2pt} & \hfill \emoji{incoming-envelope} email@email.com \\
% \raisebox{-0.05\height}\faGithub\ username \ | \
% \raisebox{-0.00\height}\faLinkedin\ username \ | \ \raisebox{-0.05\height}\faGlobe \ mysite.com  & \hfill \emoji{calling} number
% \end{tabularx}

\begin{tabularx}{\linewidth}{@{} C @{}}
\Huge{Alix Sirven-Viénot} \\[7.5pt] 
Doctorante en première année de thèse en traitement automatique des langues (Ingénierie des connaissances et Intelligence artificielle) \\
\href{https://github.com/ASVienot}{\raisebox{-0.05\height}\faGithub\ ASVienot} \ $|$ \ 
%\href{https://linkedin.com/in/alix-sirven-viénot-0238371b5}{\raisebox{-0.05\height}\faLinkedin\ username} \ $|$ \ 
% \href{https://mysite.com}{\raisebox{-0.05\height}\faGlobe \ mysite.com} \ $|$ \ 
\href{mailto:43004196@parisnanterre.fr}{\raisebox{-0.05\height}\faEnvelope \ 43004196@parisnanterre.fr} \
%\href{tel:+000000000000}{\raisebox{-0.05\height}\faMobile \ +00.00.000.000} \\
\end{tabularx}

%----------------------------------------------------------------------------------------
% EXPERIENCE SECTIONS
%----------------------------------------------------------------------------------------

%Interests/ Keywords/ Summary
\section{Sujet de thèse}
Doctorante en traitement automatique des langues (spécialité Ingénierie des connaissances et Intelligence artificielle),
 je travaille sur la caractérisation des anomalies dans le spectre des radiofréquences, en mobilisant les ontologies et le TAL pour prendre en compte l'incertitude. 
Ma thèse est réalisée sous la direction de \textbf{Pr.Delphine Battistelli} (MODYCO) et \textbf{Pr.Valentina Dragos} (ONERA) et financée en CIFRE par l'entreprise ATDI.
 %This CV can also be automatically complied and published using GitHub Actions. For details, \href{https://github.com/jitinnair1/autoCV}{click here}.

%----------------------------------------------------------------------------------------
%	EDUCATION
%----------------------------------------------------------------------------------------
\section{Formation}
\begin{tabularx}{\linewidth}{@{}p{2.3cm} X@{}}
2025 - 2029 & Doctorat \hfill \normalsize {Université Paris Nanterre, ONERA}\\

2023 - 2025 & Master Traitement automatique des langues \hfill {Paris Nanterre, Inalco, Sorbonne Nouvelle} \\ 

2022 - 2023 & MA in Applied Linguistics and Intercultural Studies \hfill {Maynooth University - Ireland} \\ 

%2019 - 2022 & Licence en Langues, Littératures et Civilisations Anglophones \hfill {Institut Catholique de Paris} \\
2019 - 2022 & Licence de langue, culture, histoire anglophone \hfill {Institut Catholique de Paris} \\


%2021 & Class 10th Some Board \hfill  (Grades) \\
\end{tabularx}


%----------------------------------------------------------------------------------------
%	SKILLS
%----------------------------------------------------------------------------------------
\section{Compétences Informatiques}
\begin{tabularx}{\linewidth}{@{}l X@{}}
Langages &  \normalsize{Python, Bash, R, LaTeX, SQL, SPARQL, OWL, RDF}\\
Ontologies &  \normalsize{Protégé, RDFlib, DOLCE, BFO}\\
Apprentissage Automatique  &  \normalsize{ScikitLearn, Pytorch, TensorFlow, HuggingFace, Pandas, Polars, Numpy}\\ 
Traitement de données &  \normalsize{Spacy, Trankit, BeautifulSoup, TreeTagger, Matplotlib}\\
Autres &  \normalsize{Git, Docker, Jupyter, VSCode, Neo4j, Linux, Windows}\\ 
\end{tabularx}

%Experience
\section{Experiences Professionelles}

\begin{jobshort}{CDI - ATDI}{Mai 2026 - present}
Contrat de travail à durée indéterminée en tant que doctorant chercheur en traitement automatique des langues pour un projet de recherche en collaboration avec l'ONERA et l'université Paris Nanterre. 
\end{jobshort}\

\begin{jobshort}{Stage - EDF R\&D Palaiseau}{Mars 2025 - Septembre 2025}
Stage de 6 mois en tant que chercheur en traitement automatique des langues pour un projet de Reconnaissance d'entités spécifiques de l'industrie de l'énergie dans des documents techniques.
\end{jobshort}\

\begin{jobshort}{CDD - EDF R\&D Palaiseau}{Septembre 2025 - Novembre 2025}
Contrat à durée déterminée afin de terminer le projet de recherche initié lors du stage de 6 mois. 
\end{jobshort}\

\begin{jobshort}{Stage - CNRS - MODYCO, Paris}{Mai 2024 - Août 2024}
Stage de 3 mois sur l'analyse et l'annotation de la notion de prise en charge dans des messages à contenu haineux issus des réseaux sociaux et de livres scolaires.
\end{jobshort}

%Projects
\section{Projects}

\begin{tabularx}{\linewidth}{ @{}l r@{} }
\textbf{Master Thesis 2023 - Applied Linguistics and Intercultural Studies} & \hfill \href{https://github.com/ASVienot/MasterThesis_Anglicisms_in_French_Medias}{Rapport} \\[3.75pt]
\multicolumn{2}{@{}X@{}}{Étude quantitative sur l’utilisation des anglicismes dans les médias Français, comparaison de leurs usages à la télévision et dans les nouveaux médias avec la plateforme YouTube.}  \\
\end{tabularx}


\begin{tabularx}{\linewidth}{ @{}l r@{} }
\textbf{Ingénierie des connaissances -- Ontologie (Master 2025)}
& \hfill \href{https://github.com/ASVienot/ontologie_masterTAL}{GitHub} \\[3.75pt]
\multicolumn{2}{@{}X@{}}{Création automatisée d'une ontologie représentant le Master TAL à partir d'une base de données et de questionnaires en ligne.} \\
\end{tabularx}

\begin{tabularx}{\linewidth}{ @{}l r@{} }
\textbf{Reconnaissance d'entités nommées -- Réseaux de neuronnes (Master 2025)}
& \hfill \href{https://github.com/KehinaleK/NER-for-Elvish-Quenya}{GitHub} \\[3.75pt]
\multicolumn{2}{@{}X@{}}{Création d'une interface permettant l'utilisation de modèles NER entrainés sur une langue fictive.} \\
\end{tabularx}

%----------------------------------------------------------------------------------------
%	PUBLICATIONS
%----------------------------------------------------------------------------------------
%\section{Publications}
%\begin{refsection}[citations.bib]
%\nocite{*}
%\printbibliography[heading=none]
%\end{refsection}

\begin{minipage}[t]{0.49\linewidth}
\section*{Expériences associatives}
\begin{itemize}[leftmargin=*, nosep, itemsep=3pt]
    \item Membre du conseil de faculté, représentant les élèves en situation de handicap, 2020-2022
    \item Bénévolat chez Rempart, préservation du patrimoine, stage de menuiserie, 2018
    \item Cheffe scout marin à Aulnay-sous-Bois, 2019-2022
\end{itemize}
\end{minipage}\hfill
\begin{minipage}[t]{0.49\linewidth}
\section*{Centres d'intérêt}
\begin{itemize}[leftmargin=*, nosep, itemsep=3pt]
    \item Art moderne
    \item Théâtre, escrime artistique
    \item Rugby, kayak, escrime artistique, vélo
    \item Loisirs créatifs (bricolage, felting, couture...)
\end{itemize}
\end{minipage}


\vfill
%\center{\footnotesize Last updated: \today}

\end{document}
